\section{Заключение}

В ходе выполнения лабораторной работы были освоены базовые возможности языка Verilog HDL на примере построения цифрового устройства из нескольких взаимодействующих модулей.

Конструкция \texttt{generate} позволила удобно масштабировать дешифратор «1 в 2» до параметризируемого модуля \texttt{decoder\_tree} с настраиваемой разрядностью входного кода. 

Для синхронного сохранения результата разработан последовательный элемент
\texttt{result\_keeper}, фиксирующий входной вектор по положительному
фронту сигнала \texttt{clk} и поддерживающий асинхронный сброс по сигналу
\texttt{rst}. Тем самым была реализована передача результата вычислений
между модулями по сигналу синхронизации, что и являлось одной из целей
работы.

Все разработанные модули объединены в тестовом окружении и проверены в симуляторе, полученные результаты полностью совпадают с ожидаемыми. Таким образом, цель работы была \textit{достигнута} в полном объёме.